\session{4}{10/17}{1限}{2}{企業活動と情報システムの全体像}

\goals{
  \item バリューチェーンの枠組みで企業活動を分け、各活動を支える情報システムを対応づけられる。
  \item 基幹系（ERP）・顧客系（CRM）・供給系（SCM）の役割と関係を説明できる。
  \item クラウド化（SaaS）の意味と、各システムにAI機能が組み込まれている現状を理解する。
}

\guidesec{解説}

\topic{バリューチェーンで企業活動を見る}
マイケル・ポーターは、企業の活動を「価値を生み出す一連の活動（価値連鎖）」として捉える枠組みを示しました。
製品やサービスが顧客に届くまでの流れに直接関わる\textbf{主活動}と、それを支える\textbf{支援活動}に分けます。
\par\noindent\begin{gtabh}{colspec={Q[l,wd=22mm] X[l] X[l]}}
  区分 & 活動 & 情報システムの例 \\
  主活動 & 購買物流、製造、出荷物流、販売・マーケティング、サービス & 調達システム、生産管理、倉庫管理（WMS）、POS・EC、コールセンター \\
  支援活動 & 全般管理（企画・財務）、人事・労務管理、技術開発、調達 & 会計システム、人事システム、設計支援（CAD）、電子購買 \\
\end{gtabh}
情報システムを学ぶとき、この枠組みは「どの活動の、どの情報を、何のために」扱うシステムなのかを整理する地図になります。
活動と活動の「つなぎ目」で情報が途切れると、在庫の過不足や納期遅れが生じます。情報システムの重要な役割は、
このつなぎ目を情報でつなぐことです。

\topic{基幹系・顧客系・供給系}
企業の情報システムは、扱う対象によって次の3つに大別されます。
\begin{itemize}
  \item \textbf{ERP（Enterprise Resource Planning、統合基幹業務システム）}　会計・人事・生産・在庫・販売などの基幹業務のデータを
        1つのデータベースで統合管理します。部門ごとに別々だったシステムを統合し、全社の状況を一貫した数字で把握できるようにします。
        導入は大規模になりがちで、業務をシステムに合わせるか、システムを業務に合わせるか（カスタマイズ）が常に論点になります。
  \item \textbf{CRM（Customer Relationship Management、顧客関係管理）}　顧客情報、購買履歴、問い合わせ、営業活動の記録を管理し、
        顧客との関係を深めるために使います。マーケティング、営業、サポートが同じ顧客情報を見られることが価値です。
  \item \textbf{SCM（Supply Chain Management、サプライチェーン管理）}　調達から生産・物流・販売に至る供給の流れを、
        自社だけでなく取引先も含めて最適化します。需要予測、在庫の最適化、納期管理が中心です。
\end{itemize}
これらに加えて、経営層や管理者の意思決定を支える\textbf{情報系システム}（BI：ビジネスインテリジェンス、データウェアハウス）があり、
基幹系のデータを集約して分析・可視化します（第8〜11回で扱います）。

\topic{クラウド化の意味}
かつて情報システムは自社でサーバーを持ち（オンプレミス）、ソフトウェアを購入して運用するものでした。
現在は、ネットワーク越しにサービスとして利用する\textbf{クラウド}が主流です。提供の形態によって、
インフラ（サーバー）を借りる IaaS、開発環境を借りる PaaS、ソフトウェアそのものを使う SaaS に分かれます。
経営の視点で重要なのは SaaS です。
\par\noindent\begin{gtabh}{colspec={Q[l,wd=24mm] X[l] X[l]}}
  観点 & オンプレミス & クラウド（SaaS） \\
  初期費用 & 大きい（設備・ライセンス） & 小さい（月額・年額の利用料） \\
  導入までの時間 & 長い & 短い \\
  機能の更新 & 自社で対応 & 提供側が継続的に更新（AI機能もここで追加される） \\
  カスタマイズ & 自由度が高い & 制約がある \\
  データの所在 & 自社内 & 事業者のサーバー（契約と設定の確認が必要） \\
  運用の負担 & 自社の情シス部門 & 提供側が担う部分が大きい \\
\end{gtabh}

\topic{AI機能が組み込まれている現状}
2023年以降、主要な業務システムには生成AIの機能が標準で組み込まれるようになりました。
オフィスソフトでの文書・メールの下書きや会議の要約、CRM での顧客対応文の生成や商談の要約、ERP での異常値の検知や
問い合わせへの自然言語での回答、SCM での需要予測の高度化などです。多くは SaaS として提供され、
「自社でAIを開発する」のではなく「使っているシステムにAIが入ってくる」形で普及しています。
経営の課題は、AIを導入するかどうかではなく、すでに入っているAI機能を何の目的で、どの業務で、どんな管理のもとで使うかを決めることに移っています。

\topic{キーワード}
\par\noindent\begin{keywords}{}
  バリューチェーン & 企業活動を主活動と支援活動に分け、どこで価値が生まれるかを捉える枠組み（ポーター）。 \\
  ERP & 会計・人事・生産・在庫などの基幹業務を1つのデータベースで統合管理するシステム。 \\
  CRM & 顧客情報・購買履歴・問い合わせを管理し、顧客との関係を深めるシステム。 \\
  SCM & 調達から販売までの供給の流れを、取引先も含めて最適化するしくみ。 \\
  BI & 基幹系のデータを集約し、分析・可視化して意思決定を支えるシステム。 \\
  SaaS & ソフトウェアをネットワーク越しにサービスとして利用する形態。提供側が機能を更新する。 \\
\end{keywords}

\guidesec{演習課題}

\exercise{1}{AI演習：身近な企業の情報システム構成をAIで推定し整理する}
\instr{よく利用する企業（コンビニ、ファストフード、通販サイト、鉄道など）を1社選び、次のプロンプトで情報システムの構成を推定させてください。出力を、バリューチェーンの活動ごとに表に整理し、「推定の根拠が示されているか」「公開情報で確かめられるか」を各行に書き添えてください。}
\begin{promptbox}[情報システム構成の推定]
あなたは企業の情報システムに詳しいITコンサルタントです。「（企業名）」の事業を、ポーターのバリューチェーンの主活動・支援活動に分け、各活動でどのような情報システム（ERP・CRM・SCM・POS・ECなど）が使われていると推定されるかを、推定の根拠（公開されている事例や業界の一般的な構成）とともに表にしてください。確実でない部分は「推定」と明記してください。
\end{promptbox}

\exercise{2}{3つのシステムの「つなぎ目」}
\instr{ある製造業で、ERP・CRM・SCM が互いに連携していないとき、どのような問題が起きるかを3つ挙げてください。それぞれ「どの情報が」「どこからどこへ」伝わらないために起きる問題かを書きます。}

\exercise{3}{SaaS を選ぶときの確認事項}
\instr{従業員50人の企業が、営業部門向けに SaaS 型の CRM を導入しようとしています。経営者として契約前に確認すべきことを、費用・機能・データ・運用の4つの観点から2つずつ挙げてください。}

\guidesec{模範解答}

\begin{answer}{演習1}
例：コンビニエンスストアチェーン
\par\noindent\begin{gtabh}{colspec={Q[l,wd=22mm] X[l] X[l,wd=46mm]}}
  活動 & 推定される情報システム & 根拠の確認 \\
  購買物流・調達 & 本部の発注・仕入システム、取引先との EDI & 業界で一般的。各社の IR 資料や採用情報で確認できる。 \\
  店舗運営（販売） & POS、店舗の発注端末、電子マネー・アプリ決済の基盤 & POS 活用は公開資料に多数あり、確かめられる。 \\
  出荷物流 & 配送センターの在庫・配送管理（温度帯別） & 推定。物流子会社の情報で一部確認できる。 \\
  マーケティング & 会員アプリ・ポイントのデータ基盤（CRM に相当） & アプリの存在は確認できる。内部構成は推定。 \\
  全般管理 & 会計・人事の ERP、BI による店舗別の分析 & 推定。ERP の製品名を挙げてきた場合は要確認。 \\
\end{gtabh}
気づき：AIは「業界の一般的な構成」としてもっともらしい表を作るが、特定企業の製品名や導入年は確かめられないことが多い。
「推定」と明記させる条件を付けても、断定的に書いてくる部分があるので、根拠の列を自分で埋めることが検証になる。
\end{answer}

\begin{answer}{演習2}
\begin{enumerate}[label=\arabic*.,leftmargin=2em]
  \item 営業が CRM に記録した大口の受注見込みが SCM に伝わらず、原材料の手配が遅れて納期を守れない（受注見込み情報：営業→調達・生産）。
  \item SCM の在庫情報が ERP の会計に反映されるのが月末の手入力のため、月中の在庫評価額が分からず資金繰りの判断が遅れる（在庫情報：物流→財務）。
  \item 顧客からの品質クレームが CRM に記録されるだけで生産側に届かず、同じ不良が繰り返される（クレーム情報：サポート→製造）。
\end{enumerate}
いずれも「情報は存在するのに、必要な部門に必要なタイミングで届かない」問題です。システム連携は技術の問題であると同時に、
部門間の業務の設計の問題でもあります。
\end{answer}

\begin{answer}{演習3}
\begin{itemize}
  \item 費用：ユーザー数に応じた月額の総額（3年分で試算）。追加機能・サポート・データ移行の費用。
  \item 機能：自社の営業プロセス（見積→受注→請求）に必要な機能があるか。既存の会計システムやメールと連携できるか。
  \item データ：顧客データの保管場所（国内か国外か）と、解約時にデータを取り出せるか。入力した顧客情報がAI機能の学習に使われないか。
  \item 運用：誰が設定と権限管理を担当するか。障害時の連絡体制と稼働率の保証（SLA）。
\end{itemize}
\end{answer}

\guidesec{出席確認クイズの解説}
講義サイトの出席確認ページ（第4回）の5問です。
% 出席確認クイズ 第04回　企業活動と情報システムの全体像
%   attendance/attendance04.html から生成（解説は本ガイド用に加筆）。

\quizq{1}{ERP(統合基幹業務システム)の説明として適切なものはどれでしょうか?}%
  {顧客との関係を管理するシステム}%
  {\correct{会計・人事・在庫などの基幹業務を統合して管理するシステム}}%
  {電子メールを配信するシステム}%
  {社内向けのSNS}%
  {正解は (b)。ERP（Enterprise Resource Planning）は、会計・人事・生産・在庫などの基幹業務のデータを1つのデータベースで統合管理するシステムです。(a) は CRM の説明です。部門ごとにばらばらだったデータが統合されることで、全社の状況を一貫した数字で把握できるようになります。}

\quizq{2}{CRMが主に扱う対象はどれでしょうか?}%
  {原材料の調達}%
  {\correct{顧客との関係(購買履歴や問い合わせなど)}}%
  {建物の管理}%
  {従業員の給与}%
  {正解は (b)。CRM（Customer Relationship Management）は、顧客情報・購買履歴・問い合わせなどを管理し、顧客との関係を深めるためのしくみです。(a) の調達は SCM、(d) の給与は ERP の人事給与モジュールが担います。}

\quizq{3}{SCMの目的として適切なものはどれでしょうか?}%
  {社内文書を保管する}%
  {\correct{調達から生産・物流・販売までの供給の流れを最適化する}}%
  {ソフトウェアを開発する}%
  {顧客の満足度を調査する}%
  {正解は (b)。SCM（Supply Chain Management）は、調達から生産・物流・販売に至る供給の流れ全体を、在庫や納期の面で最適化します。(a)(c)(d) は文書管理・開発・調査で、供給連鎖の最適化とは別の活動です。}

\quizq{4}{「バリューチェーン(価値連鎖)」の考え方を提唱したのは誰でしょうか?}%
  {スティーブ・ジョブズ}%
  {\correct{マイケル・ポーター}}%
  {アダム・スミス}%
  {ピーター・ドラッカー}%
  {正解は (b)。バリューチェーン（価値連鎖）はマイケル・ポーターが『競争優位の戦略』（1985年）で示した枠組みで、企業活動を5つの主活動と4つの支援活動に分けて、どこで価値が生まれるかを分析します。ドラッカーはマネジメント論、スミスは分業論で知られ、ジョブズは経営者です。}

\quizq{5}{クラウド(SaaS)型の情報システムにAI機能が組み込まれることの利点として適切なものはどれでしょうか?}%
  {従業員の教育が不要になる}%
  {\correct{自社で開発しなくても最新のAI機能を利用できる}}%
  {データが社外に出ることが絶対にない}%
  {ネットワークが不要になる}%
  {正解は (b)。SaaS 型のサービスでは提供側が機能を更新するため、自社で開発しなくても最新のAI機能を使えます。一方で (c) は誤りで、データは事業者のサーバーに置かれるため契約と設定の確認が必要です。(a) 教育は引き続き必要で、(d) ネットワークはむしろ不可欠です。}

